% =========================================================
% CHAPITRE 2 — CONCEPTION ET ÉTUDE DES BESOINS
% =========================================================

\section{Introduction}

Après avoir présenté le contexte général du projet, sa problématique et la
solution proposée, ce chapitre est consacré à l'étude des besoins et à la
conception du système \textit{3awedlou}.

L'objectif de cette étape est de définir les fonctionnalités attendues,
les contraintes auxquelles le système doit répondre ainsi que
l'architecture générale permettant d'assurer la communication entre la
machine, les systèmes embarqués, le serveur et les applications clientes.

La conception est organisée autour de plusieurs aspects complémentaires :
l'analyse des besoins fonctionnels et non fonctionnels, l'étude de
l'enchaînement des données, la définition de l'architecture matérielle et
logicielle, l'étude des protocoles de communication IoT ainsi que la
modélisation UML du système.

\section{Besoins fonctionnels}

Les besoins fonctionnels décrivent les services et les fonctionnalités que
le système doit fournir à ses différents utilisateurs et composants.

Dans le cadre du projet \textit{3awedlou}, les principales fonctionnalités
identifiées sont les suivantes.

\subsection{Gestion des utilisateurs}

Le système doit permettre aux utilisateurs autorisés d'accéder à
l'application et d'utiliser les fonctionnalités correspondant à leurs
droits d'accès. L'authentification permet notamment de sécuriser l'accès
aux informations relatives aux machines et aux productions.

\subsection{Gestion des machines}

L'application doit permettre de gérer les machines connectées au système.
Pour chaque machine, les informations relatives à son identification et à
son état doivent être accessibles.

Le système doit notamment permettre de connaître l'état courant d'une
machine, par exemple :

\begin{itemize}
    \item en attente ;
    \item en cours de chauffage ;
    \item en extrusion ;
    \item en pause ;
    \item en erreur ;
    \item hors ligne.
\end{itemize}

\subsection{Acquisition des données}

La machine doit être capable de mesurer et de transmettre différentes
données relatives à son fonctionnement. Ces données peuvent notamment
concerner la température, l'état du moteur, l'état général de la machine
et les informations relatives à une session de production.

\subsection{Supervision de la machine}

Le système doit permettre à l'utilisateur de consulter l'état de la
machine et les données collectées. La supervision permet ainsi de suivre
le fonctionnement du système sans avoir à intervenir directement sur la
machine.

\subsection{Gestion des sessions de fonctionnement}

Une session doit permettre d'identifier une période de fonctionnement de
la machine. Le système doit pouvoir enregistrer les informations
associées au démarrage, au fonctionnement et à l'arrêt d'une session.

\subsection{Suivi de la production}

Le système doit permettre d'enregistrer les informations relatives à la
production de filament. Ces informations peuvent notamment comprendre
l'identification d'un lot, les paramètres de production et les données
associées au processus de transformation.

\subsection{Communication temps réel}

La solution doit permettre la transmission des données entre la machine et
le système logiciel avec un délai suffisamment faible pour assurer une
supervision efficace.

Cette communication repose sur une architecture IoT utilisant le protocole
MQTT, qui sera étudié plus en détail dans les sections suivantes.

\section{Besoins non fonctionnels}

Les besoins non fonctionnels définissent les contraintes et les propriétés
qualitatives que le système doit respecter au-delà de ses fonctionnalités
principales.

\subsection{Fiabilité}

Le système doit fonctionner de manière stable pendant les différentes
phases de fonctionnement de la machine. Les échanges de données doivent
également être suffisamment fiables pour éviter la perte d'informations
importantes.

\subsection{Sécurité}

L'accès aux fonctionnalités de l'application et aux données du système
doit être protégé. L'authentification des utilisateurs ainsi que la
sécurisation des communications constituent donc des éléments importants
de l'architecture.

\subsection{Performance}

Le système doit pouvoir traiter les données provenant de la machine sans
introduire de délais importants dans la supervision. Les communications
doivent donc rester suffisamment rapides pour permettre un suivi proche du
temps réel.

\subsection{Évolutivité}

L'architecture doit permettre l'ajout de nouvelles machines, de nouveaux
capteurs ou de nouvelles fonctionnalités sans nécessiter une
reconception complète du système.

\subsection{Maintenabilité}

La séparation entre la partie matérielle, les systèmes embarqués, le
backend et les applications clientes doit faciliter la maintenance et
l'évolution du projet.

\subsection{Ergonomie}

Les interfaces de supervision doivent présenter les informations de
manière claire et compréhensible. L'utilisateur doit pouvoir identifier
rapidement l'état de la machine et les principales données de production.

\section{Schéma d'enchaînement des données}

Le fonctionnement global du système repose sur une succession d'échanges
de données entre la machine, les systèmes embarqués, le serveur et les
applications clientes.

La machine constitue la source principale des données. Les capteurs
permettent d'acquérir les informations relatives à son fonctionnement,
tandis que l'Arduino Mega assure une partie du contrôle local des
actionneurs et du processus d'extrusion.

L'ESP32 joue le rôle d'interface de communication entre la partie
embarquée et l'infrastructure réseau. Les données provenant de la machine
sont transmises par l'intermédiaire du protocole MQTT vers le broker IoT.

Le backend récupère ensuite les messages nécessaires, traite les données
et les enregistre dans la base de données. Les applications web et mobile
peuvent alors récupérer et afficher ces informations à l'utilisateur.

Le cheminement général des données peut être résumé comme suit :

\begin{center}
\textbf{
Machine $\rightarrow$ Arduino Mega $\rightarrow$ ESP32
$\rightarrow$ MQTT Broker $\rightarrow$ Backend
$\rightarrow$ Applications Web / Mobile
}
\end{center}

Dans le sens inverse, les commandes provenant des applications peuvent
suivre le chemin opposé afin d'être transmises jusqu'à la machine.

\section{Architecture du système}

L'architecture de \textit{3awedlou} est conçue selon une approche
distribuée permettant de séparer les responsabilités des différents
composants.

Elle peut être divisée en plusieurs niveaux :

\begin{enumerate}
    \item \textbf{Couche physique :} machine de recyclage, capteurs,
    moteur, système de chauffage et autres actionneurs.
    
    \item \textbf{Couche embarquée :} Arduino Mega et ESP32 assurant
    respectivement le contrôle local et la communication réseau.
    
    \item \textbf{Couche IoT :} broker MQTT assurant l'échange de
    messages entre les équipements connectés et les services logiciels.
    
    \item \textbf{Couche serveur :} backend chargé de l'authentification,
    de la gestion des machines, du traitement des données et de leur
    stockage.
    
    \item \textbf{Couche application :} applications web et mobile
    permettant la supervision et l'exploitation des données.
\end{enumerate}

Cette séparation permet de réduire le couplage entre les différentes
parties du système et facilite son évolution.

\section{Protocoles IoT}

La communication constitue un élément essentiel d'une architecture IoT.
Le choix du protocole doit tenir compte des caractéristiques du système,
notamment la quantité de données échangées, les ressources disponibles,
la latence souhaitée et la fiabilité des communications.

Plusieurs protocoles peuvent être envisagés pour assurer la communication
entre les équipements et les services logiciels.

\subsection{MQTT}

MQTT (\textit{Message Queuing Telemetry Transport}) est un protocole de
communication léger basé sur un modèle \textit{publish/subscribe}. Les
équipements publient des messages sur des sujets (\textit{topics}) et les
clients intéressés s'abonnent aux sujets correspondants.

Cette architecture permet de découpler les producteurs et les
consommateurs de données.

\subsection{HTTP/HTTPS}

HTTP est largement utilisé pour les communications entre applications
clientes et serveurs. Il repose principalement sur un modèle
requête/réponse et convient particulièrement aux API web.

Dans une architecture IoT, HTTP peut être utilisé pour certaines
opérations nécessitant une communication avec le backend, mais il est
moins adapté qu'un protocole orienté messages pour certaines
communications temps réel entre équipements.

\subsection{WebSocket}

WebSocket permet d'établir une communication persistante entre un client
et un serveur. Il est particulièrement adapté aux applications nécessitant
des mises à jour en temps réel de l'interface utilisateur.

Il peut donc constituer une solution intéressante pour la communication
entre le backend et les interfaces web ou mobiles.

\section{Choix du protocole}

Après comparaison des différentes solutions, MQTT a été retenu comme
protocole principal de communication IoT du projet.

Ce choix est principalement motivé par son modèle
\textit{publish/subscribe}, qui permet de découpler la machine du backend
et des applications clientes.

Le broker MQTT joue le rôle d'intermédiaire central. La machine peut ainsi
publier ses données sans avoir à connaître directement les applications
qui les utiliseront. De la même manière, les applications ou le backend
peuvent publier des commandes sur des sujets auxquels l'ESP32 est abonné.

Cette architecture présente également l'avantage de faciliter
l'intégration de plusieurs machines dans le futur.

MQTT est donc utilisé comme mécanisme principal pour les échanges de
données IoT, tandis que les technologies web classiques peuvent être
utilisées pour les interactions entre les applications clientes et le
backend.

\section{Diagramme UML}

La modélisation UML permet de représenter le comportement du système et
les interactions entre ses différents acteurs et composants.

Dans le cadre du projet \textit{3awedlou}, plusieurs diagrammes UML sont
utilisés afin de représenter les fonctionnalités principales, le
déroulement des processus et les échanges entre les composants.

\subsection{Diagramme de cas d'utilisation}

Le diagramme de cas d'utilisation permet d'identifier les acteurs
interagissant avec le système ainsi que les principales fonctionnalités
auxquelles ils ont accès.

Dans le système \textit{3awedlou}, l'utilisateur peut notamment consulter
l'état d'une machine, suivre les données de fonctionnement et consulter
les informations relatives à la production.

Le système embarqué constitue également un élément important de
l'architecture puisqu'il fournit les données issues de la machine et peut
recevoir des commandes provenant du système logiciel.

% Insérer ici le diagramme de cas d'utilisation.
%
% Exemple :
% \begin{figure}[H]
%     \centering
%     \includegraphics[width=0.9\textwidth]{figures/use-case.png}
%     \caption{Diagramme de cas d'utilisation}
%     \label{fig:use-case}
% \end{figure}

\subsection{Diagramme d'activité}

Le diagramme d'activité permet de représenter l'enchaînement des
principales étapes du fonctionnement du système.

Il peut notamment représenter le processus allant du démarrage de la
machine jusqu'à la production du filament, en passant par les phases de
chauffage, de vérification des conditions de fonctionnement et
d'extrusion.

% Insérer ici le diagramme d'activité.
%
% \begin{figure}[H]
%     \centering
%     \includegraphics[width=0.9\textwidth]{figures/activity-diagram.png}
%     \caption{Diagramme d'activité du système}
%     \label{fig:activity}
% \end{figure}

\subsection{Diagramme de séquence}

Le diagramme de séquence permet de représenter les échanges de messages
entre les différents composants du système au cours d'un scénario donné.

Dans le cas de \textit{3awedlou}, un scénario de supervision peut par
exemple suivre le cheminement d'une donnée de télémétrie depuis la machine
jusqu'à l'application.

La donnée est d'abord acquise par le système embarqué, puis transmise à
l'ESP32. Celui-ci publie le message sur le broker MQTT. Le backend reçoit
le message, traite les informations et les enregistre si nécessaire. Les
applications peuvent ensuite exploiter ces données afin de mettre à jour
l'interface de supervision.

% Insérer ici le diagramme de séquence.
%
% \begin{figure}[H]
%     \centering
%     \includegraphics[width=0.95\textwidth]{figures/sequence-diagram.png}
%     \caption{Diagramme de séquence de transmission des données}
%     \label{fig:sequence}
% \end{figure}

\section{Conclusion}

Ce chapitre a permis d'identifier les besoins fonctionnels et non
fonctionnels du système \textit{3awedlou} et de définir son architecture
générale.

L'étude a également permis de décrire le cheminement des données entre la
machine, les systèmes embarqués, l'infrastructure IoT, le backend et les
applications clientes. Plusieurs protocoles ont été étudiés et MQTT a été
retenu comme protocole principal pour les communications IoT en raison de
son modèle publish/subscribe et de son adéquation avec les besoins du
projet.

Enfin, les diagrammes UML permettent de compléter la conception en
représentant les fonctionnalités du système, le déroulement de ses
processus et les échanges entre ses différents composants.

Le chapitre suivant sera consacré à la réalisation du système, notamment
à la mise en œuvre de la partie matérielle, des systèmes embarqués, de la
communication IoT, du backend et des interfaces de supervision.
